\begin{theorem}[Hasse-function lift]\label{mod:finitefield-hassefunctionlift}
Let $K/k$ be a finite extension of finite fields, let
$\psi\in\operatorname{FiniteAddChar}(k)$ be nontrivial, and let
$\phi\in\operatorname{HasseFunction}(k,\psi)$ as in
Definition~\ref{mod:finitefield-hassefunction}.  For $x\in K$, put
\[
 p_x=\operatorname{minpoly}_k(x),\qquad
 h_x=[K:k(x)],\qquad
 s_x=-\operatorname{nextCoeff}(p_x),\qquad
 e_x=\operatorname{secondCoeff}(p_x),
\]
and define the unordered-pair correction
\[
 \omega_{K/k}(x)=h_xe_x+\binom{h_x}{2}s_x^2.
\]
This is the second coefficient of the repeated minimal polynomial
$p_x^{h_x}$.  Equivalently, it is the elementary symmetric sum
\[
 \sum_{\{\sigma,\tau\}}\sigma(x)\tau(x)
\]
over unordered pairs of distinct $k$-embeddings of $K$ into an
algebraic closure, with every pair occurring once.  In particular, it
is not the ordered off-diagonal sum.

Define
\[
 \phi_{K/k}(x)
 =\phi\!\left(\Tr_{K/k}x\right)
   \psi\!\left(-\omega_{K/k}(x)\right).
\]
Writing
\(\psi_{K/k}(z)=\psi(\Tr_{K/k}z)\), the correction satisfies
\[
 \omega_{K/k}(x+y)
 =\omega_{K/k}(x)+\omega_{K/k}(y)
   +\Tr_{K/k}(x)\Tr_{K/k}(y)-\Tr_{K/k}(xy),
\]
and hence the lifted function is nowhere zero and satisfies the Hasse
functional equation
\[
 \phi_{K/k}(x+y)
 =\phi_{K/k}(x)\phi_{K/k}(y)\psi_{K/k}(xy).
\]
Thus \(\phi_{K/k}\in\operatorname{HasseFunction}(K,\psi_{K/k})\).
Moreover, the exact lift identity is
\[
 \boxed{
 -\sum_{x\in K}\phi_{K/k}(x)
 =\left(-\sum_{x\in k}\phi(x)\right)^{[K:k]} .}
\]
The trace is directed from $K$ down to $k$, the correction has the
displayed negative sign inside $\psi$, and the exponent is the extension
degree $[K:k]$.

More explicitly, if $p=X^d+a_1X^{d-1}+a_2X^{d-2}+\cdots$ is monic and
$h\geq0$, then
\[
 \operatorname{nextCoeff}(p^h)=h a_1,
 \qquad
 \operatorname{secondCoeff}(p^h)
 =h a_2+\binom h2a_1^2.
\]
The second formula is integral and characteristic-independent.  Indeed,
if the $d$ conjugates are each repeated $h$ times, pairs with different
conjugate values contribute $h^2e$, while equal-value pairs contribute
$\binom h2\sum_i\xi_i^2$.  Using $\sum_i\xi_i^2=s^2-2e$ gives
$he+\binom h2s^2$ without dividing by two.

For the polynomial weight
\[
 w(P)=\phi\!\left(-\operatorname{nextCoeff}(P)\right)
       \psi\!\left(-\operatorname{secondCoeff}(P)\right)
\]
from Definition~\ref{mod:finitefield-polynomialweights}, its exact-degree
sums satisfy
\[
 A_1(w)=\sum_{x\in k}\phi(x),
 \qquad
 A_n(w)=0\quad(n\geq2).
\]
The first equality includes the convention that the second coefficient
of a linear polynomial is zero.  The second uses the explicit hypothesis
$\psi\ne1$: translation of the freely varying second coefficient
multiplies the sum by a nontrivial value of $\psi$.

For $r=[K:k]$, partitioning $K$ by minimal polynomial identifies the
closed-point power sum of Theorem~\ref{mod:finitefield-eulerproduct} with
the lifted sum:
\[
 B_r(w)=\sum_{x\in K}\phi_{K/k}(x).
\]
A monic irreducible polynomial of degree $d\mid r$ has exactly $d$ roots
in $K$, and for each such root $x$ one has $h_x=r/d$ and
\[
 \phi_{K/k}(x)=w(p_x)^{h_x}.
\]
Thus the factor $d$ and the power $r/d$ in $B_r$ are both retained.
Applying the reusable identity $-B_r(w)=(-A_1(w))^r$ gives the boxed
formula.  Its proof uses $A_0(w)=1$ but does not divide by a weight or by
the Hasse sum, and it assumes no nonvanishing beyond the nowhere-zero
values built into $\operatorname{HasseFunction}$ and the stated
nontriviality of $\psi$.
\LeanFile{LanglandsFirstMainLemma/FiniteField/HasseFunctionLift.lean}
\lean{LanglandsFirstMainLemma.hasseEmbeddingMultiplicity}
\lean{LanglandsFirstMainLemma.hassePairCorrection}
\lean{LanglandsFirstMainLemma.hasseFunctionLiftValue}
\lean{LanglandsFirstMainLemma.monic_nextCoeff_pow}
\lean{LanglandsFirstMainLemma.monic_secondCoeff_pow}
\lean{LanglandsFirstMainLemma.lmul_charpoly_eq_minpoly_pow}
\lean{LanglandsFirstMainLemma.hassePairCorrection_eq_secondCoeff_minpoly_pow}
\lean{LanglandsFirstMainLemma.hassePairCorrection_eq_secondCoeff_charpoly}
\lean{LanglandsFirstMainLemma.hassePairCorrection_eq_elementarySymmetric_two}
\lean{LanglandsFirstMainLemma.Finset.esymm_two_map_add}
\lean{LanglandsFirstMainLemma.hasseGaloisPairSum}
\lean{LanglandsFirstMainLemma.hasseGaloisPairSum_add}
\lean{LanglandsFirstMainLemma.hasseGaloisLinearFactorProduct_eq_minpoly_pow}
\lean{LanglandsFirstMainLemma.algebraMap_hassePairCorrection_eq_hasseGaloisPairSum}
\lean{LanglandsFirstMainLemma.hassePairCorrection_add}
\lean{LanglandsFirstMainLemma.hasseFunctionLiftValue_map_add}
\lean{LanglandsFirstMainLemma.hasseFunctionTraceLift}
\lean{LanglandsFirstMainLemma.hasseFunctionTraceLift_apply}
\lean{LanglandsFirstMainLemma.residueTrace_eq_neg_nextCoeff_minpoly_pow}
\lean{LanglandsFirstMainLemma.hassePolynomialWeight_pow}
\lean{LanglandsFirstMainLemma.hasse_chooseTwo_correction}
\lean{LanglandsFirstMainLemma.hasseFunctionLiftValue_eq_minpolyWeight_pow}
\lean{LanglandsFirstMainLemma.degreeSum_hassePolynomialWeight_one}
\lean{LanglandsFirstMainLemma.degreeSum_hassePolynomialWeight_eq_zero}
\lean{LanglandsFirstMainLemma.hasseMinpolyFiberEquivAlgHom}
\lean{LanglandsFirstMainLemma.hasseMinpolyFiber_card}
\lean{LanglandsFirstMainLemma.hasseClosedPointIndexEquiv}
\lean{LanglandsFirstMainLemma.hasse_degree_div_eq_embeddingMultiplicity}
\lean{LanglandsFirstMainLemma.hasseClosedPointPowerSum_eq_liftSum}
\lean{LanglandsFirstMainLemma.hasseFunctionLift}
\leanok
\SourceText{Theorem thm:Hasse-lift, with the unordered-pair convention for distinct embeddings.}
\uses{mod:finitefield-hassefunction,mod:finitefield-eulerproduct,mod:localfield-extension}
\FormalizationContract
\end{theorem}
\begin{proof}
The coefficient formulas follow by induction from the first- and
second-coefficient product identities.  Multiplicativity of $w$ then
shows directly that the Hasse value of a repeated closed point is
$w(p)^h$; equivalently, the factor
$\psi(\binom h2s^2)$ supplied by the Hasse functional equation cancels
the same term in $\psi(-\omega)$.

To verify the functional equation of the lift itself, identify
$\omega_{K/k}$ with the second elementary symmetric function of the
multiset of Galois conjugates.  This is the unordered-pair sum: its
polarization has cross term
$\Tr(x)\Tr(y)-\Tr(xy)$, with no factor of two.  Substitution into
$\psi(-\omega(x+y))$ cancels the factor
$\psi(\Tr(x)\Tr(y))$ contributed by the Hasse equation for $\phi$ and
leaves exactly $\psi(\Tr(xy))=\psi_{K/k}(xy)$.  The product defining the
lift is nowhere zero because both character values and all values of
$\phi$ are nonzero.

Linear monic polynomials are parametrized by $X+a$, and negation
reindexes their weights to the Hasse sum.  In degrees at least two,
adding $cX^{n-2}$ permutes monic polynomials of degree $n$, fixes the
first coefficient, and adds $c$ to the second coefficient.  Choosing
$c$ with $\psi(-c)\ne1$ forces the degree-$n$ sum to vanish.

For the closed-point comparison, roots with fixed minimal polynomial
$p$ are identified with $k$-algebra homomorphisms $k[X]/(p)\to K$.
Finite-field embedding counts give exactly $d$ such roots when
$d=\deg p$ divides $[K:k]$.  The tower degree formula gives
$h=[K:k]/d$.  Reindexing the sum over $K$ by these fibers therefore
gives $B_{[K:k]}(w)$, and
Theorem~\ref{mod:finitefield-eulerproduct} supplies the final signs and
extension-degree power.
\end{proof}
